\section{总结与延伸}

本讲从“什么是通用视觉表示”出发，经过 VTAB、BiT、ViT、Scaling ViT 与 MLP-Mixer，形成一套比架构记忆更持久的方法论：

\begin{enumerate}
  \item \textbf{先定义适配目标}：通用表示应降低新任务的样本、优化与规格成本；
  \item \textbf{用任务族而非单榜测量}：natural、specialized、structured tasks 揭示不同泛化维度；
  \item \textbf{区分 upstream 与 downstream}：预训练表示和适配器容量必须分别控制；
  \item \textbf{无标签不等于无先验}：SSL 把设计转移到 proxy objective 与 augmentation；
  \item \textbf{规模是联合变量}：数据、容量、训练时长与 regularization 必须匹配；
  \item \textbf{去重是科学有效性要求}：近重复会把记忆伪装成迁移；
  \item \textbf{ViT 是弱化视觉先验的实验}：patch、position 与 class token 仍构成明确设计；
  \item \textbf{数据决定归纳偏置的价值区间}：强先验提高低数据效率，灵活模型可能在大数据追上；
  \item \textbf{Compute-normalized 仍不等于系统效率}：FLOPs、latency、memory 与 utilization 要分别报告；
  \item \textbf{Shape 和 schedule 是 scaling 的一部分}：大模型不能用小模型的短训练制度评价；
  \item \textbf{Scaling law 是局部经验模型}：可用于规划，但必须报告拟合区间与转折风险；
  \item \textbf{Mixer 的意义是反事实比较}：attention 并非唯一机制，完整 recipe 才是因果对象。
\end{enumerate}

\begin{importantbox}{课程作业：设计一个可信的视觉表示评估}
选择一个预训练模型与三个差异明显的 downstream tasks。定义 upstream provenance、去重方法、linear probe 与 full fine-tuning 两套 protocol；报告每组任务的 mean/worst-case、10-shot/full-data、calibration、training/inference FLOPs 与真实 latency。最后写出至少三个“结果不能证明什么”，并说明若数据预算扩大十倍，你预期哪种模型族的曲线斜率会变化。
\end{importantbox}

\begin{warningbox}{最终自检：不要把五种成功混在一起}
\textbf{Upstream fit}、\textbf{in-distribution accuracy}、\textbf{few-shot transfer}、\textbf{robustness} 与 \textbf{deployment efficiency} 是五种不同成功。论文或系统若只证明其中一项，结论就应停在那一项；不能用 ImageNet top-1 替代通用表示，也不能用理论 FLOPs 替代真实服务成本。
\end{warningbox}

\subsection{拓展阅读}

{\small
\begin{itemize}[nosep,leftmargin=*]
  \item \href{https://arxiv.org/abs/1910.04867}{\textbf{VTAB}}：多任务视觉表示适配 benchmark 与 task-family 设计。
  \item \href{https://arxiv.org/abs/1912.11370}{\textbf{Big Transfer (BiT)}}：大规模监督预训练、迁移 recipe 与 robustness。
  \item \href{https://arxiv.org/abs/2010.11929}{\textbf{An Image is Worth 16x16 Words}}：ViT 架构、数据 scaling 与 transfer。
  \item \href{https://arxiv.org/abs/2106.10270}{\textbf{How to Train Your ViT?}}：augmentation、regularization、dataset size 与 training recipe。
  \item \href{https://arxiv.org/abs/2106.04560}{\textbf{Scaling Vision Transformers}}：model shape、schedule、few-shot 与 scaling laws。
  \item \href{https://arxiv.org/abs/2105.01601}{\textbf{MLP-Mixer}}：token/channel mixing 与 data-scale comparison。
\end{itemize}
}

\end{document}
